% GENERATED FILE. Edit canonical lessons, then run npm run generate:books.
% canonical-source-hash: fnv1a64:f785d6c3ef11f538
% canonical-lessons: SA-C380-matra, SA-C380-parimanam, SA-C380-madhyamah, SA-C380-yathestam, SA-C380-vistarayati

\chapter{How Much}
\label{ch:sa-how-much}
\hlchaptermodality{380}

\begin{chapteropening}
\textbf{By the end of this chapter:} \emph{I can talk about amounts, sizes and degrees.}

Complete the last lesson of chapter 380: I can talk about amounts, sizes and degrees.
\end{chapteropening}

\section[\texorpdfstring{\sk{मात्रा} (mātrā)}{मात्रा (mātrā)}]{\sk{मात्रा} (mātrā) — a quantity; a measure}

\label{lesson:SA-C380-matra}



\begin{quote}
Before the new one: say the Sanskrit for ample, then the Sanskrit for increases (something).
\end{quote}

\subsection*{You'll want to know: \sk{मात्रा}}
\textbf{\sk{मात्रा}} — \emph{mātrā} — \textquotedblleft{}a quantity; a measure\textquotedblright{}.

\begin{quote}
\sk{यदि} \sk{औषधस्य} \sk{मात्रा} \sk{अधिका}, \sk{तर्हि} \sk{हानिः} \sk{भविष्यति।}
\end{quote}

\begin{quote}
\emph{yadi auṣadhasya mātrā adhikā, tarhi hāniḥ bhaviṣyati.}
\end{quote}

\begin{quote}
If the dose of medicine is too large, it will do harm.
\end{quote}

A feminine noun.

\begin{cousinweb}
\textbf{\sk{मात्रा}} (\emph{mātrā}) comes from the root √mā, \textquotedblleft{}to measure\textquotedblright{}; MW gives \emph{mātrā} \textquotedblleft{}measure, quantity, size, duration, degree\textquotedblright{} (MW).
\end{cousinweb}

\begin{grammarlens}[title={What you've built}]
One more word for saying what happens if.
\end{grammarlens}

\subsection*{Your turn}
\begin{itemize}
\raggedright
  \item \emph{Say it:} \emph{mātrā}
  \item \emph{Say it:} \emph{mātrā}, once more
  \item \emph{Say it:} \emph{yadi auṣadhasya mātrā adhikā, tarhi hāniḥ bhaviṣyati.} — If the dose of medicine is too large, it will do harm.
  \item \emph{Recall:} say the Sanskrit for ample, then the Sanskrit for increases (something), then say \emph{mātrā} again
\end{itemize}

\begin{tcolorbox}[breakable,colback=teal!4,colframe=teal!35!black,title={Before you move on}]
What does \sk{मात्रा} mean? (\textquotedblleft{}A quantity; a measure\textquotedblright{}.) Say it once more.
\end{tcolorbox}
\section[\texorpdfstring{\sk{परिमाणम्} (parimāṇam)}{परिमाणम् (parimāṇam)}]{\sk{परिमाणम्} (parimāṇam) — an amount; a size}

\label{lesson:SA-C380-parimanam}



\begin{quote}
Before the new one: say the Sanskrit for increases (something), then the Sanskrit for a quantity.
\end{quote}

\subsection*{You'll want to know: \sk{परिमाणम्}}
\textbf{\sk{परिमाणम्}} — \emph{parimāṇam} — \textquotedblleft{}an amount; a size\textquotedblright{}.

\begin{quote}
\sk{यदि} \sk{जलस्य} \sk{परिमाणं} \sk{न्यूनम्}, \sk{तर्हि} \sk{वृक्षाः} \sk{शुष्यन्ति।}
\end{quote}

\begin{quote}
\emph{yadi jalasya parimāṇaṃ nyūnam, tarhi vṛkṣāḥ śuṣyanti.}
\end{quote}

\begin{quote}
If the amount of water is low, the trees dry up.
\end{quote}

A neuter noun.

\begin{cousinweb}
\textbf{\sk{परिमाणम्}} (\emph{parimāṇam}) is built on \emph{pari}, \textquotedblleft{}around\textquotedblright{}, and the root √mā, \textquotedblleft{}to measure\textquotedblright{}; MW gives \emph{pari-māṇa} \textquotedblleft{}measure of any kind, size, weight, number\textquotedblright{} (MW).
\end{cousinweb}

\begin{grammarlens}[title={What you've built}]
One more word for saying what happens if.
\end{grammarlens}

\subsection*{Your turn}
\begin{itemize}
\raggedright
  \item \emph{Say it:} \emph{parimāṇam}
  \item \emph{Say it:} \emph{parimāṇam}, once more
  \item \emph{Say it:} \emph{yadi jalasya parimāṇaṃ nyūnam, tarhi vṛkṣāḥ śuṣyanti.} — If the amount of water is low, the trees dry up.
  \item \emph{Recall:} say the Sanskrit for increases (something), then the Sanskrit for a quantity, then say \emph{parimāṇam} again
\end{itemize}

\begin{tcolorbox}[breakable,colback=teal!4,colframe=teal!35!black,title={Before you move on}]
What does \sk{परिमाणम्} mean? (\textquotedblleft{}An amount; a size\textquotedblright{}.) Say it once more.
\end{tcolorbox}
\section[\texorpdfstring{\sk{मध्यमः} (madhyamaḥ)}{मध्यमः (madhyamaḥ)}]{\sk{मध्यमः} (madhyamaḥ) — moderate; middling}

\label{lesson:SA-C380-madhyamah}



\begin{quote}
Before the new one: say the Sanskrit for a quantity, then the Sanskrit for an amount.
\end{quote}

\subsection*{You'll want to know: \sk{मध्यमः}}
\textbf{\sk{मध्यमः}} — \emph{madhyamaḥ} — \textquotedblleft{}moderate; middling\textquotedblright{}.

\begin{quote}
\sk{यदि} \sk{व्यायामः} \sk{मध्यमः}, \sk{तर्हि} \sk{देहः} \sk{स्वस्थः} \sk{भवति।}
\end{quote}

\begin{quote}
\emph{yadi vyāyāmaḥ madhyamaḥ, tarhi dehaḥ svasthaḥ bhavati.}
\end{quote}

\begin{quote}
If exercise is moderate, the body stays healthy.
\end{quote}

\begin{cousinweb}
\textbf{\sk{मध्यमः}} (\emph{madhyamaḥ}) is built on \emph{madhya}, \textquotedblleft{}middle\textquotedblright{}, with the superlative ending \emph{-ma}; MW gives \emph{madhyama} \textquotedblleft{}middle, middlemost, intermediate\textquotedblright{} (MW).
\end{cousinweb}

\begin{grammarlens}[title={What you've built}]
One more word for saying what happens if.
\end{grammarlens}

\subsection*{Your turn}
\begin{itemize}
\raggedright
  \item \emph{Say it:} \emph{madhyamaḥ}
  \item \emph{Say it:} \emph{madhyamaḥ}, once more
  \item \emph{Say it:} \emph{yadi vyāyāmaḥ madhyamaḥ, tarhi dehaḥ svasthaḥ bhavati.} — If exercise is moderate, the body stays healthy.
  \item \emph{Recall:} say the Sanskrit for a quantity, then the Sanskrit for an amount, then say \emph{madhyamaḥ} again
\end{itemize}

\begin{tcolorbox}[breakable,colback=teal!4,colframe=teal!35!black,title={Before you move on}]
What does \sk{मध्यमः} mean? (\textquotedblleft{}Moderate; middling\textquotedblright{}.) Say it once more.
\end{tcolorbox}
\section[\texorpdfstring{\sk{यथेष्टम्} (yatheṣṭam)}{यथेष्टम् (yatheṣṭam)}]{\sk{यथेष्टम्} (yatheṣṭam) — as much as one likes}

\label{lesson:SA-C380-yathestam}



\begin{quote}
Before the new one: say the Sanskrit for an amount, then the Sanskrit for moderate.
\end{quote}

\subsection*{You'll want to know: \sk{यथेष्टम्}}
\textbf{\sk{यथेष्टम्}} — \emph{yatheṣṭam} — \textquotedblleft{}as much as one likes\textquotedblright{}.

\begin{quote}
\sk{यदि} \sk{भवान्} \sk{क्षुधितः}, \sk{तर्हि} \sk{यथेष्टं} \sk{खादतु।}
\end{quote}

\begin{quote}
\emph{yadi bhavān kṣudhitaḥ, tarhi yatheṣṭaṃ khādatu.}
\end{quote}

\begin{quote}
If you are hungry, eat as much as you like.
\end{quote}

\begin{cousinweb}
\textbf{\sk{यथेष्टम्}} (\emph{yatheṣṭam}) is built from \emph{yathā}, \textquotedblleft{}as\textquotedblright{}, and \emph{iṣṭa}, \textquotedblleft{}wished\textquotedblright{}; MW gives \emph{yatheṣṭam} \textquotedblleft{}according to wish, at pleasure\textquotedblright{} (MW).
\end{cousinweb}

\begin{grammarlens}[title={What you've built}]
One more word for saying what happens if.
\end{grammarlens}

\subsection*{Your turn}
\begin{itemize}
\raggedright
  \item \emph{Say it:} \emph{yatheṣṭam}
  \item \emph{Say it:} \emph{yatheṣṭam}, once more
  \item \emph{Say it:} \emph{yadi bhavān kṣudhitaḥ, tarhi yatheṣṭaṃ khādatu.} — If you are hungry, eat as much as you like.
  \item \emph{Recall:} say the Sanskrit for an amount, then the Sanskrit for moderate, then say \emph{yatheṣṭam} again
\end{itemize}

\begin{tcolorbox}[breakable,colback=teal!4,colframe=teal!35!black,title={Before you move on}]
What does \sk{यथेष्टम्} mean? (\textquotedblleft{}As much as one likes\textquotedblright{}.) Say it once more.
\end{tcolorbox}
\section[\texorpdfstring{\sk{विस्तारयति} (vistārayati)}{विस्तारयति (vistārayati)}]{\sk{विस्तारयति} (vistārayati) — expands; extends}

\label{lesson:SA-C380-vistarayati}



\begin{quote}
Before the new one: say the Sanskrit for moderate, then the Sanskrit for as much as one likes.
\end{quote}

\subsection*{You'll want to know: \sk{विस्तारयति}}
\textbf{\sk{विस्तारयति}} — \emph{vistārayati} — \textquotedblleft{}expands; extends\textquotedblright{}.

\begin{quote}
\sk{यदि} \sk{वणिक्} \sk{व्यापारं} \sk{विस्तारयति}, \sk{तर्हि} \sk{लाभः} \sk{वर्धते।}
\end{quote}

\begin{quote}
\emph{yadi vaṇik vyāpāraṃ vistārayati, tarhi lābhaḥ vardhate.}
\end{quote}

\begin{quote}
If the merchant expands his trade, his profit grows.
\end{quote}

Past (he …): \sk{सः} \sk{व्यापारं} \sk{विस्तारितवान्}, \emph{saḥ vyāpāraṃ vistāritavān}.

\begin{cousinweb}
\textbf{\sk{विस्तारयति}} (\emph{vistārayati}) is built on \emph{vi} and the root √stṛ, \textquotedblleft{}to spread\textquotedblright{}, in its causative form; MW gives \emph{vi-stārita} \textquotedblleft{}spread, extended\textquotedblright{} (MW).
\end{cousinweb}

\begin{grammarlens}[title={What you've built}]
One more word for saying what happens if.
\end{grammarlens}

\subsection*{Your turn}
\begin{itemize}
\raggedright
  \item \emph{Say it:} \emph{vistārayati}
  \item \emph{Say it:} \emph{vistārayati}, once more
  \item \emph{Say it:} \emph{yadi vaṇik vyāpāraṃ vistārayati, tarhi lābhaḥ vardhate.} — If the merchant expands his trade, his profit grows.
  \item \emph{Recall:} say the Sanskrit for moderate, then the Sanskrit for as much as one likes, then say \emph{vistārayati} again
\end{itemize}

\begin{tcolorbox}[breakable,colback=teal!4,colframe=teal!35!black,title={Before you move on}]
What does \sk{विस्तारयति} mean? (\textquotedblleft{}Expands; extends\textquotedblright{}.) Say it once more.
\end{tcolorbox}
